For a field \(\F\), vectors \(a \in \F^n\) and \(b \in \F^m\), the tensor product \(a \otimes b \in \F^{n \times m}\)
is the matrix defined by
\[
  (a \otimes b)_{i,j} := a_i b_j \quad \forall i \in [n],\, j \in [m].
\]
We write \(\mathrm{flat}(a \otimes b) \in \F^{nm}\) for the row-concatenation of \(a \otimes b\).

For our purposes, we want a test that checks whether $\mathrm{flat}(a \otimes a) = b$

\begin{definition}\label{TENSORQ}\lean{TENSORQ}\leanok
For a field \(\F\) and \(n \in \N\), define the tensor test language
\[
  \mathrm{TENSORQ}(\F, n) :=
    \bigl\{ (a, b) \in \F^n \times \F^{n^2} \mid b = \mathrm{flat}(a \otimes a) \bigr\}.
\]
\end{definition}

\begin{definition}\label{TENSORQ.verifier}\lean{TENSORQ.verifier}\uses{LPCPVerifier,TENSORQ}\leanok
% Lecture 08 presents this as a subcheck inside the QESAT verifier, where the
% queried proof oracles are already named \(a\) and \(b\). The standalone Lean
% language has explicit input \((a,b)\), so its verifier also checks that the
% proof's two linear parts agree with that input before running the tensor
% multiplicativity check.
Given \((a,b) \in \F^n \times \F^{n^2}\), the LPCP verifier
\(V_{\mathrm{TENSORQ}(\F,n)}^{\langle \pi,\cdot\rangle}(a,b)\) uses a proof
\(\pi=\alpha\parallel\beta \in \F^{n+n^2}\), intended to be the linear encoding of the public input
\((a,b)\). Since the proof may be dishonest, the verifier checks consistency with \((a,b)\) on the sampled
queries before checking the tensor relation:
\begin{enumerate}
  % \item Sample \(s,t \leftarrow \F^n\).
  % \item Query the \(a\)-part of \(\pi\) at \(s\) and at \(t\), obtaining \(\langle a, s \rangle\) and \(\langle a, t \rangle\).
  % \item Query the \(b\)-part of \(\pi\) at \(\mathrm{flat}(s\otimes t)\), obtaining \(\langle b, \mathrm{flat}(s \otimes t) \rangle\).
  % \item Accept if and only if \(\langle b, \mathrm{flat}(s \otimes t) \rangle = \langle a, s \rangle \cdot \langle a, t \rangle\).
  \item Sample \(s,t \leftarrow \F^n\).
  \item Query the \(\alpha\)-part of \(\pi\) at \(s\) and at \(t\), obtaining \(y_s\) and \(y_t\).
  \item Query the \(\beta\)-part of \(\pi\) at \(\mathrm{flat}(s\otimes t)\), obtaining \(y_{s,t}\).
  \item Accept if and only if
  \[
      y_s=\langle a,s\rangle,\qquad
      y_t=\langle a,t\rangle,\qquad
      y_{s,t}=\langle b,\mathrm{flat}(s\otimes t)\rangle,\qquad
      y_{s,t}=y_s y_t.
    \]
\end{enumerate}
\end{definition}

\begin{definition}\label{TENSORQ.selfVerifier}\lean{TENSORQ.selfVerifier}\uses{OracleComp,randOracleSpec}\leanok
In settings where \(a\) and \(b\) are supplied only by the proof oracle, the tensor self-check uses the same three
queries but has no public input to compare against. Given oracle access to
\(\pi=\alpha\parallel\beta\in\F^{n+n^2}\), sample \(s,t\leftarrow\F^n\), query
\[
  y_s=\langle\alpha,s\rangle,\qquad
  y_t=\langle\alpha,t\rangle,\qquad
  y_{s,t}=\langle\beta,\mathrm{flat}(s\otimes t)\rangle,
\]
and accept if and only if \(y_{s,t}=y_s y_t\). This is the tensor component used inside the QESAT verifier.
\end{definition}

\begin{lemma}\label{TENSORQ.tensor_check_complete}\lean{TENSORQ.tensor_check_complete}\leanok
For every \(a,s,t\in\F^n\),
\[
  \langle \mathrm{flat}(a\otimes a),\mathrm{flat}(s\otimes t)\rangle
  = \langle a,s\rangle \langle a,t\rangle.
\]
\end{lemma}
\begin{proof}\leanok
\begin{align*}
  \langle \mathrm{flat}(a\otimes a),\mathrm{flat}(s\otimes t)\rangle
    &= \sum_{i,j} a_i a_j s_i t_j
     = \left(\sum_i a_i s_i\right)\left(\sum_j a_j t_j\right)
     = \langle a,s\rangle \langle a,t\rangle.
\end{align*}
\end{proof}

\begin{lemma}\label{TENSORQ.verifier_soundness_after_sampling}\lean{TENSORQ.verifier_soundness_after_sampling}\uses{TENSORQ.verifier}\leanok
Let \(\F\) be a finite field. If \(b \neq \mathrm{flat}(a\otimes a)\), then for every proof
\(\pi\in \F^{n+n^2}\),
\[
  \Pr_{s,t}\!\left[
    V_{\mathrm{TENSORQ}(\F,n)}^{\langle\pi,\cdot\rangle}(a,b)=1
  \right]
  \leq \frac{2|\F|-1}{|\F|^2}.
\]
\end{lemma}
\begin{proof}\uses{TENSORQ.verifier,LINEQ.linear_form_uniform_prob_mul_card_le_one}\leanok
Suppose \(b \neq \mathrm{flat}(a \otimes a)\), i.e., there exists \(i^*, j^* \in [n]\) such that
\(b_{i^* j^*} \neq a_{i^*} \cdot a_{j^*}\).
For each \(i \in [n]\) define the linear polynomial \(p_i(t) := \sum_{j \in [n]} (b_{ij} - a_i a_j)\, t_j\).
For fixed \(s,t\), if the verifier accepts, then its consistency checks force
\(y_s=\langle a,s\rangle\), \(y_t=\langle a,t\rangle\), and
\(y_{s,t}=\langle b,\mathrm{flat}(s\otimes t)\rangle\). Thus acceptance can only occur when
\[
  \langle b, \mathrm{flat}(s \otimes t) \rangle - \langle a,s\rangle\langle a,t\rangle
  = \sum_{i \in [n]}  \sum_{j \in [n]} (b_{ij} - a_i a_j)\, s_i t_j
  = \sum_{i \in [n]} p_i(t)\, s_i.
\]
Immediate application of Polynomial Identity Testing yields a lower bound of $1 - \tfrac{2}{|\F|}$ but we can do better.
Since \(b_{i^* j^*} \neq a_{i^*} a_{j^*}\), the polynomial \(p_{i^*}\) is nonzero, so by the
Schwartz--Zippel lemma applied to \(t\),
\[
  \Pr_{t}\!\left[p_{i^*}(t) \neq 0\right] \geq 1 - \tfrac{1}{|\F|}.
\]
Conditioned on \(p_{i^*}(t) \neq 0\), the expression \(\sum_i p_i(t)\,s_i\) is a nonzero linear polynomial
in \(s\), so by the Schwartz--Zippel lemma applied to \(s\),
\[
  \Pr_{s \leftarrow \F^n}\!\left[\sum_i p_i(t)\,s_i \neq 0 \mid p_{i^*}(t) \neq 0\right]
  \geq 1 - \tfrac{1}{|\F|}.
\]
Hence
\[
  \Pr_{s,t}\!\left[ V_{\mathrm{TENSORQ}(\F,n)}^{\langle\pi,\cdot\rangle}(a,b)=0 \right]
  \geq \left(1 - \tfrac{1}{|\F|}\right)^{\!2},
\]
and therefore
\[
  \Pr_{s,t}\!\left[ V_{\mathrm{TENSORQ}(\F,n)}^{\langle\pi,\cdot\rangle}(a,b)=1 \right]
  \leq 1 - \left(1 - \tfrac{1}{|\F|}\right)^{\!2}
  = \frac{2|\F|-1}{|\F|^2}.
\]
\end{proof}

\begin{lemma}\label{TENSORQ.selfVerifier_soundness_after_sampling}\lean{TENSORQ.selfVerifier_soundness_after_sampling}\uses{TENSORQ.selfVerifier}\leanok
Let \(\F\) be a finite field and let \(\pi=\alpha\parallel\beta\in\F^{n+n^2}\). If
\((\alpha,\beta)\notin\mathrm{TENSORQ}(\F,n)\), then the tensor self-check accepts with probability at most
\[
  \frac{2|\F|-1}{|\F|^2}.
\]
\end{lemma}
\begin{proof}\uses{TENSORQ.verifier_soundness_after_sampling}\leanok
Apply \ref{TENSORQ.verifier_soundness_after_sampling} to the public pair \((\alpha,\beta)\). The self-check is
exactly the multiplicativity part of that verifier, with \((\alpha,\beta)\) read from the proof oracle itself.
\end{proof}


\begin{theorem}\label{TENSORQ_LPCP}\uses{TENSORQ, LPCP}\lean{TENSORQ_LPCP}\leanok
For every finite field \(\F\) and \(n \in \N\),
\[
  \mathrm{TENSORQ}(\F, n) \in
  \LPCP\!\left[
    \varepsilon_c = 0,
    \varepsilon_s = \frac{2|\F|-1}{|\F|^2},
    \Sigma = \F,
    \ell = n + n^2,
    q = 3,
    r = 2n
  \right].
\]
\end{theorem}
\begin{proof}\uses{TENSORQ.verifier,TENSORQ.tensor_check_complete,TENSORQ.verifier_soundness_after_sampling}\leanok
Use the verifier of \ref{TENSORQ.verifier}. If
\((a,b)\in\mathrm{TENSORQ}(\F,n)\), the honest proof \(\pi=(a,b)\) passes the
three input-consistency checks identically, and the multiplicativity check is
\ref{TENSORQ.tensor_check_complete}. Soundness is
\ref{TENSORQ.verifier_soundness_after_sampling}. The verifier makes three proof
queries and uses two random vectors in \(\F^n\), which gives the displayed query
and randomness bounds.
\end{proof}

\begin{corollary}\label{TENSORQ_LPCP_Zmod2}\uses{TENSORQ, LPCP}\leanok
\[
  \mathrm{TENSORQ}(\F_2, n) \in
  \LPCP\!\left[
    \varepsilon_c = 0,
    \varepsilon_s = 3 / 4,
    \Sigma = \F_2,
    \ell = n + n^2,
    q = 3,
    r = 2 n
  \right].
\]
\end{corollary}
\begin{proof}\uses{TENSORQ_LPCP}\lean{TENSORQ_LPCP_Zmod2}\leanok
Specialize \ref{TENSORQ_LPCP} to \(\F_2\), where \(|\F_2|=2\).
\end{proof}
